\section{Change of algebraically closed field}\label{sec:field}

The preceding arguments use complex topology.  We now explain
both transfers needed for Theorem~\ref{thm:main}: a whole
countable hypothetical program is transferred to $\mathbb C$
for the termination argument, whereas existence is transferred
one finite diagram at a time.  These are different assertions;
existence of one terminating complex program would not suffice
for the first of them.

\begin{lemma}\label{lem:numerical-extension}
Let $k_0\subset K$ be algebraically closed fields of
characteristic zero and let $X_0$ be a projective variety over
$k_0$.  Base extension identifies their real numerical divisor
and curve spaces and their closed effective curve cones:
\[
 N^1(X_0)_{\mathbb R}\simeq N^1(X_{0,K})_{\mathbb R},
 \qquad
 N_1(X_0)_{\mathbb R}\simeq N_1(X_{0,K})_{\mathbb R},
 \qquad
 \overline{\mathrm{NE}}(X_0)
    =\overline{\mathrm{NE}}(X_{0,K}).
\]
In fact every numerical class of an integral curve on the
extension is represented by an integral curve over $k_0$.
For a projective morphism between projective varieties, these
identifications commute with pullback of numerical divisor
classes and preserve the relative Picard number.
\end{lemma}

\begin{proof}
Injectivity on numerical divisors follows by testing against
curves defined over $k_0$.  For surjectivity, a line bundle
$\mathcal L$ on $X_{0,K}$ spreads to a line bundle on
$X_0\times S$, where $S$ is an integral finite-type
$k_0$-scheme and the given map $\operatorname{Spec}K\to S$
is dominant.  Shrink $S$ if necessary and choose a closed point
$s\in S(k_0)$.  The line bundles $\mathcal L$ and
$(\mathcal L_s)_K$ have the same numerical class on $X_{0,K}$.
Indeed they are fibres of the same line-bundle family on the
fixed variety $X_{0,K}$.  Since $k_0$ is algebraically closed, $S$
is geometrically integral; hence the parameter scheme $S_K$ is
connected.  The family's degree on any fixed integral curve in
$X_{0,K}$ is constant.  This last assertion is the constancy of degree
in a line-bundle family on a proper curve, or follows from
constancy of its Euler characteristic.  Thus every numerical
divisor class comes from $k_0$.

Similarly, an integral curve on $X_{0,K}$ spreads to a flat
projective family of geometrically integral curves in
$X_0\times S$ after shrinking an integral finite-type parameter
space.  Geometric integrality is available since the given
extension field is algebraically closed.  Specialize to a
$k_0$-point of this open parameter space.  Degrees against every
line bundle on $X_0$ are unchanged by flatness, so the resulting
integral curve represents the same numerical class.  The
already established statement for numerical divisors makes
these precisely all necessary degree tests.  Conversely, an
integral curve over an algebraically closed field stays
integral on algebraically closed extension.  The effective
curve classes, their positive spans, and their closures
therefore identify.  The usual nondegenerate numerical
pairings give the asserted identification of curve spaces.

Pullback commutes with base extension.  For a projective morphism
$f_0:X_0\to Z_0$ with projective target, choose an ample line
bundle $H$ on $Z_0$.  An integral curve is contracted precisely
when its degree against $f_0^*H$ is zero.  If a curve on the
extension is contracted, its integral numerical representative
over $k_0$ has the same zero degree and is also contracted.
Conversely, every contracted integral curve over $k_0$ stays
integral and contracted on extension.  The identifications
therefore preserve the span of contracted curve classes in
$N_1(X_0)_{\mathbb R}$.  The relative numerical divisor space
is dual to this span, so its dimension, the relative Picard
number, is preserved.
\end{proof}

\begin{proposition}\label{prop:field-transfer}
Let $k$ be any algebraically closed field of characteristic
zero.  A countable program of the kind in
Definition~\ref{def:program}, including all its prescribed ray
choices, descends to a common countable algebraically closed
subfield of $k$.  After embedding that subfield in $\mathbb C$,
it gives a program with the same step types and relative signs,
the same boundary coefficients, and the same singularity
properties.  In particular an infinite birational program
over $k$ would give an infinite birational program over
$\mathbb C$.
\end{proposition}

\begin{proof}
At each index the varieties, maps, boundaries, chosen rational
canonical divisors, and their compatibility involve only
finitely many coefficients.  Include equations for projective
embeddings, ample line bundles, relative ample multiples of
the two adjoints, the Cartier multiples defining their
$\mathbb Q$-Cartier property, and the relevant bundle
isomorphisms.  Include a log resolution at each index, with
its SNC divisors and crepant pullback coefficients.  There are
only countably many such data.  They are therefore all defined
over one countably generated field over $\mathbb Q$.
Its algebraic closure inside $k$ is a countable algebraically
closed field $k_0$ over which the whole diagram is defined.
It admits an embedding into $\mathbb C$.

These algebraically closed extensions preserve normality,
integrality, dimension, projectivity, birationality, smallness,
and whether a morphism is an isomorphism.  Relative ampleness
is preserved and is detected by faithful field extension.
The chosen canonical divisors and Cartier data preserve the
actual adjoints and the strict-transform or pushforward
boundary identities.  The descended log resolutions are
smooth with the same SNC data.  Their crepant coefficients
preserve lc and klt wherever asserted.  For lc, the calculation
in Lemma~\ref{lem:snc-valuations} uses only the identity
$a(v)=\sum_jv(T_j)a_j+e(v)$ with $a_j\geq0$ and $e(v)\geq0$;
the coefficients $a_j$ may vanish.  No finiteness assertion
requiring their positivity is used.

Lemma~\ref{lem:numerical-extension}, applied at every index,
preserves extremality, negativity, and the relative Picard
number of each chosen contraction.  An ample polarization on
its target identifies the contracted integral curves by degree
zero against its pullback, so these curve classes are preserved
as well.  Connected fibres of the birational maps follow from
properness and normality.  For a prescribed fibre-type endpoint,
the identity $f_*\mathcal O_X=\mathcal O_Z$ is detected over $k_0$
by proper flat base change along the faithful field extension
$k_0\subset k$, and preserved by the same theorem along
$k_0\subset\mathbb C$.  This keeps the prescribed endpoint
itself.  Each nontrivial birational step remains nontrivial.
Consequently all indices of an infinite proposed program
survive under this single embedding.

Only the original countable program and the data testing its
properties are descended here.  The auxiliary terminalizations,
valuation constructions, and Hodge-theoretic objects in the
contradiction may then be constructed over $\mathbb C$.
\end{proof}

\begin{lemma}\label{lem:finite-existence-transfer}
Proposition~\ref{prop:continuation} holds over any algebraically
closed characteristic-zero field, with exactly the contraction
and mixed-step conventions stated there.
\end{lemma}

\begin{proof}
Let $(X,B)$ be a pair over $k$ and let $R$ be a negative extremal
ray.  Choose a countable algebraically closed field of
definition $k_0\subset k$ for $(X,B)$ and embed $k_0$ into
$\mathbb C$.  By Lemma~\ref{lem:numerical-extension}, $R$
identifies with a negative extremal ray on both $X_0$ and
$X_{0,\mathbb C}$.  Proposition~\ref{prop:continuation} produces
the required complex contraction and, in the birational case,
its lc positive model.

These outputs are finite-presentation data.  Spread them over
an integral finite-type $k_0$-scheme $S$ whose function field
embeds in $\mathbb C$, keeping the input equal to the fixed
product $X_0\times S$.  Include projective embeddings and
polarizations of the targets, Cartier multiples and relative
ample bundles for the adjoints, the boundary divisors and
strict-transform identities, and log resolutions with their
crepant coefficients.  After shrinking $S$, the geometric
fibres of the targets are normal and integral with the same
dimensions, and the required morphisms are surjective and
projective.  In the birational case they are birational, and
the positive map is small.  To see the latter assertions
directly, spread the isomorphisms on dense opens together with
their inverses.  The complements on the positive source can
be kept of dimension at most two; this excludes an exceptional
prime in every remaining fibre.  Smoothness and the SNC
conditions on the finitely many resolutions, and the exact
crepant pullback formulas, can likewise be retained after
shrinking.  The resulting boundaries are the prescribed
transforms and the resulting pairs are lc.

The relative signs are retained by spreading ample line-bundle data.
For example, choose a relatively very ample power of each
relevant adjoint multiple and a sufficiently positive base
twist making it generated.  Spread the resulting projective
embedding.  Its restriction preserves relative ampleness.
All these requirements concern a finite diagram and hold
on a nonempty open subset of $S$.

We check that the \emph{particular} ray is preserved, including
its relative Picard number.  Write
$f_{\mathbb C}:X_{0,\mathbb C}\to Z_{\mathbb C}$ for the complex
contraction.  Choose an ample line bundle $H$ on its target
and include it in the spread.  Its pullback is a line bundle
on the fixed product $X_0\times S$, so its numerical class is
constant on the connected parameter space, by the argument
of Lemma~\ref{lem:numerical-extension}.  Thus
\[
 \alpha=[f_{\mathbb C}^*H]\in N^1(X_0)_{\mathbb R}
\]
is the same at every specialization.

Choose an integral curve $C_0$ over $k_0$ with nonzero class in
$R$, using the integral-curve part of
Lemma~\ref{lem:numerical-extension} and a curve spanning the
complex negative ray.  Its degree against a fixed ample line
bundle on $X_0$ is positive.  For any closed point $s\in S(k_0)$
in the chosen open and every integral curve $C\subset X_0$,
\[
 \begin{aligned}
 f_s(C)\text{ is a point}
 &\quad\Longleftrightarrow\quad \alpha\cdot C=0\\
 &\quad\Longleftrightarrow\quad
 f_{\mathbb C}(C_{\mathbb C})\text{ is a point}\\
 &\quad\Longleftrightarrow\quad [C]\in R.
 \end{aligned}
\]
The first two equivalences follow from the projection formula,
ampleness of the target polarizations, and constancy of $\alpha$.
The last is the defining property of the original complex
contraction.  Hence $f_s$ contracts precisely the given ray.
Since it contracts $C_0$, the span of its contracted curve
classes is the nonzero one-dimensional span of $R$.
Its relative Picard number is therefore one.  The fixed class
$\alpha$ thus controls the contracted curves at every specialization.

For a birational output, connected fibres follow from
proper birationality and normality of the targets.  In the
fibre-type case, one may take the Stein factorization of
$f_s$.  Its intermediate target is normal because $X_0$ is
normal, is finite over $Z_s$, and has the same dimension.
The pullback of $H_s$ to that target remains ample.  Consequently
the contracted integral curves and the span of their classes
are unchanged, so the relative Picard number remains one.
Relative ampleness of $-D$ persists under this
factorization.  We now have connected fibres and the required
fibre-type endpoint.

Choose a $k_0$-point of the nonempty open parameter space;
such a point exists because $k_0$ is algebraically closed.
The preceding construction gives the desired finite diagram
over $k_0$.  Extend it to $k$.  Normality, lc singularities,
the relative signs, the exact contracted ray, the relative
Picard number, and the stipulated smallness are preserved.
This is the required output for the original chosen ray.
\end{proof}

\begin{proof}[Proof of Theorem~\ref{thm:main}]
If an allowed program over $k$ had infinitely many birational
steps, Proposition~\ref{prop:field-transfer} would transfer
that entire prescribed program to $\mathbb C$.  This contradicts
Proposition~\ref{prop:lc-lift}.  Thus every allowed program
has finitely many birational steps.

At any stage with non-nef adjoint, there is a negative
extremal ray.  This follows from the lc cone theorem, or from
its complex form using Lemma~\ref{lem:numerical-extension}.
Every selected such ray has the contraction and output of
Lemma~\ref{lem:finite-existence-transfer}.  A finite
prefix at which neither stopping condition holds therefore extends.  A maximal program
can stop only at a nef pair or at a fibre-type contraction.
The latter has connected fibres, smaller-dimensional normal
projective target, relative Picard number one, and relatively
ample negative adjoint, exactly as required for the Mori
fibre-space endpoint in Definition~\ref{def:program}.
An initially nef pair needs no operation.

These arguments prove continuation, finiteness for every sequence of
permitted choices, and the asserted alternatives for a maximal program.
\end{proof}
